\section{The cubic CM refinement at a fixed ordinary modular form}
\label{sec:cubic-cm-euler-refinement}

We retain the actual ordinary flag when referring to a type-II
point.  The statements below do not assert that every point of a
new local deformation component is ordinary.  They first treat
an attached newform whose level is prime to the coefficient prime.

\begin{lemma}[AC-K29-ray-class: the auxiliary ray-class condition]
\label{AC-K29-ray-class}
Let $K=\mathbb Q(\omega)$ and let $\chi$ be the cubic character
of $K(\sqrt[3]{269})/K$.  Its finite conductor is
$\mathfrak f=(269)$.  The ray class group of this modulus has order
\[
       h_{\mathfrak f}=\frac{269^2-1}{6}
               =12060=2^2\,3^2\,5\,67.                 \tag{CER.1}
\]
If $\ell\ge5$ splits in $K$, $\ell\ne269$, and
$\chi(\operatorname{Frob}_{\mathfrak p})\ne1$ for a prime
$\mathfrak p\mid\ell$, then $\ell\nmid h_{\mathfrak f}$.
Moreover $\chi(\operatorname{Frob}_{\mathfrak p})$ and
$\chi(\operatorname{Frob}_{\overline{\mathfrak p}})$ remain
distinct modulo every coefficient prime above $\ell$.
\end{lemma}
\begin{proof}
The prime $269\equiv2\pmod3$ is inert in $K$.  At that prime,
the Kummer equation is tamely and totally ramified of degree
three; a nontrivial character therefore has conductor exponent
one.  At a prime different from $3,269$, adjoining a cube root
of a unit gives an unramified extension, since cubing is an
automorphism of the principal units.  At $3$, the polynomial
$X^3-269$ has a root in $\mathbb Q_3$: at $X=-1$ its value
has valuation three and its derivative has valuation one, so
the strict Hensel inequality $3>2$ applies.  Thus the character
is trivial on $G_{K_3}$.  This proves the conductor assertion.

The ring $\mathbb Z[\omega]$ has class number one, and its six
units inject into $(\mathcal O_K/(269))^\times$.
The ray class exact sequence consequently gives (CER.1),
because the residue field has $269^2$ elements.  A split prime
$\ell\ge5$ cannot be $5$.  The only remaining possible divisor
of (CER.1) is $67$.  But $269\equiv1\pmod{67}$, so $269$ is
a cube in both residue fields at $67$.  Hensel's lemma, now
with derivative a unit, shows that the Kummer character is
locally trivial there.  This excludes $67$ under the stated
Frobenius condition.  Finally the two values are $\omega$ and
$\omega^{-1}$, in one order; their difference has square
equal to $-3$, and is therefore a unit above $\ell\ge5$.
\end{proof}

\begin{proposition}[AC-K29-ray: a fixed-form CM ray Euler system]
\label{AC-K29-ray}
Let $\ell\ge5$ and $K,\chi,\mathfrak p$ satisfy
AC-K29-ray-class.  Let $f$ be a cuspidal newform of weight
$2e+1$, $e\ge1$, nebentypus $\eta_K$, and level $N_f$
prime to $\ell$.  Suppose $f$ is ordinary at a chosen
coefficient prime above $\ell$.  Write $M_f$ for the
cohomological representation in the convention of
Loeffler--Zerbes (2016), and put, for $a\in\{1,2\}$,
\[
 V=M_f(e),\qquad R_a=V|_{G_K}\chi^a,
 \qquad U_a=R_a^*(1)=M_f^*|_{G_K}\chi^{-a}(1-e).
                                                            \tag{CER.2}
\]
There are ray-class Euler systems for $U_a$, obtained from
the two ordinary CM refinements of $\theta_\chi$.
They have a stable lattice and a single nonzero denominator
independent of the auxiliary squarefree ray modulus.
Their classes are Bloch--Kato locally away from $\ell$.
With $K(\mathfrak n)$ denoting the maximal abelian
$\ell$-extension of conductor dividing $\mathfrak n$, their
norm relation at every permitted new auxiliary prime is
\[
 \operatorname{cor}_{K(\mathfrak n\mathfrak q)/K(\mathfrak n)}
       z_{\mathfrak n\mathfrak q}
       =P_{\mathfrak q}([\mathfrak q])z_{\mathfrak n},
 \qquad
 P_{\mathfrak q}(X)
       =\det(1-X\operatorname{Fr}_{\mathfrak q}^{-1}\mid R_a).
                                                            \tag{CER.3}
\]
Here $\operatorname{Fr}_{\mathfrak q}$ is arithmetic
Frobenius, while $[\mathfrak q]$ is the geometrically
normalized Artin symbol, as in the source norm relation.
One may exclude a fixed finite set containing the conductors,
$\ell$, and the smoothing integer.

Under Shapiro's isomorphism the two base classes agree, up to
nonzero coefficient scalars, with the two ordinary Hida
Beilinson--Flach specializations of KLZ Definition~8.1.1 for
\[
  (f,g,j)=(f,\theta_\chi,e-1),\qquad
 M_f^*\otimes M_g^*(-j)=\operatorname{Ind}_K^{\mathbb Q}U_a.
                                                            \tag{CER.4}
\]
No irreducibility assumption on the residual representation
of $f$, no even-weight assumption, and no trivial-nebentypus
assumption is needed for this fixed-form construction.
\end{proposition}
\begin{proof}
\proofstep{1. The universal class precedes residual localization.}
Use the integral ordinary class of
Kings--Loeffler--Zerbes, arXiv:1503.02888v6,
Theorem~6.3.4, printed pp.~34--35.  It is a class on the
product of modular curves with two Iwasawa coefficient sheaves
and a cyclotomic Iwasawa factor.  In its construction the
ordinary projector and the inverse unit-root operators act
integrally; neither a residual representation nor a choice of
cuspidal Hecke branch has yet been imposed.
For different levels in the two factors use the asymmetric
pushforward of the same class.  Independence of the common
multiple follows from its first norm relation.

\proofstep{2. A single fixed first-factor quotient.}
Keep the first tame level equal to $N_f$ for every auxiliary
modulus.  Apply the moment map of degree $2e-1$, followed by
the fixed ordinary $f$-quotient.  The global ordinary control
map is KLZ Proposition~7.2.1(4), printed pp.~35--36; this
statement is for the full ordinary module and does not require
a non-Eisenstein maximal ideal.  The passage from the ordinary
stabilization to $M_f^*$ is the map in \S7.3.  In the present
positive degree $2e-1$, Proposition~7.3.1 gives an integral
isomorphism on the indicated quotient lattices.  Alternatively,
any fixed characteristic-zero quotient has bounded image in a
chosen lattice after multiplication by one fixed scalar.
This scalar is independent of the second level because the
first module and map are fixed.

We do not identify the whole localized branch with its
parabolic module.  The November 2023 correction, KLZ \S13,
printed p.~67, explicitly warns against that identification
at an Eisenstein congruence.  The full-module moment map and
the fixed classical cuspidal quotient used here are unaffected.

\proofstep{3. Only the CM factor requires integral patching.}
For the second factor take the ordinary CM family through
$\theta_{\chi^a}$, choosing either prime above $\ell$.
Its residual induced character is irreducible: the cubic
characters $\chi^a,\chi^{-a}$ are distinct in characteristic
$\ell\ne3$.  Its two ordinary residual characters at $\ell$
are distinct by AC-K29-ray-class.  Thus the integral CM
freeness and norm isomorphisms of LLZ (2015),
Propositions~5.2.3 and~5.2.5 and Corollary~5.2.6,
printed pp.~16--18, apply.  Their simultaneous ordinary
version is precisely B\"uy\"ukboduk--Lei,
arXiv:1602.07508v3, Proposition~3.7 and Corollaries~3.8--3.9,
printed pp.~17--18.  The ray-class order hypothesis there
is supplied by (CER.1).

More explicitly, the CM cohomology at modulus
$\mathfrak f\mathfrak n\mathfrak p^r$ is a free rank-two
module over the corresponding $\ell$-primary ray group ring.
The compatible integral maps $\nu_{\mathfrak f\mathfrak n,r}$
identify it with an induced rank-one character.  Tensor these
maps with the already fixed first-factor quotient and apply
Shapiro.  Finally specialize the CM weight to one and the
cyclotomic character to $e-1$.  All these specializations are
maps on integral Iwasawa modules.  The CM weight-one class is
obtained from this family specialization, not from an
algebraic moment of degree minus one.  The resulting
coefficient is exactly $U_a$ in (CER.2).  Interchanging the
two primes gives the other CM type.  Equivalently one may
construct the systems for $a=1,2$ with a fixed CM type and
then conjugate one by the nontrivial automorphism of $K$;
since $V$ extends to $G_{\mathbb Q}$, this identifies the
coefficients with $U_a$ and exchanges the auxiliary primes
in the norm relations.
Since $\ell\nmid h_{\mathfrak f}$, the maximal
$\ell$-ray extensions at $\mathfrak f\mathfrak n$ and at
$\mathfrak n$ agree.  This explains the fields in (CER.3).

Every map on the varying CM side is integral and norm
compatible.  The fixed first-factor scalar of Step~2
therefore bounds the denominators for all $\mathfrak n$ at
once.  No inverse matrix whose denominator grows with
$\mathfrak n$ has been introduced.

\proofstep{4. The degree-two norm calculation.}
The norm identities are imposed before either specialization.
They are KLZ Theorems~5.3.1 and~5.5.1, together with the
asymmetric pushforward diagrams in BL Proposition~3.10,
printed pp.~19--21.  The latter proof applies these identities
and the CM norm map, then applies the first-factor Hecke
eigencharacter.  In particular its calculation uses
$T_q(f)=a_q(f)$ and the diamond eigenvalue
$\varepsilon_f(q)$ as scalars; it nowhere replaces the latter
by one or uses that the weight is even.

One can check this directly using LLZ Theorem~3.5.1 and
Appendix~A: the two degeneracy maps give a linear expression
in the first-factor Hecke operator and the two CM characters.
After the CM norm combination, the result is
$1-\operatorname{tr}(\operatorname{Fr}_{\mathfrak q}^{-1}
\mid R_a)[\mathfrak q]
+\det(\operatorname{Fr}_{\mathfrak q}^{-1}
\mid R_a)[\mathfrak q]^2$.
The twist $j=e-1$ changes the dual Euler coefficient to
$M_f\chi^a(1+j)=R_a$, exactly as in BL (7).
The remaining split, inert, and repeated-prime cases are
the same operator identities before the first-factor
quotient, so applying that quotient preserves them.
This proves (CER.3).  The residual-image assumptions used
later in BL's Selmer and Hida-branch reciprocity arguments
are not inputs to these universal identities.

\proofstep{5. Local conditions and the base class.}
Before finite cyclotomic specialization the construction is
an Iwasawa cohomology class.  At each place away from
$\ell$, its rational specializations are in the unramified
Bloch--Kato subspace.  This is the universal-norm statement
used in Loeffler--Zerbes (2016), Proposition~2.4.4, printed
p.~10; it applies over every $K(\mathfrak n)$.

At $\mathfrak n=1$, the compatible CM identification is
the ordinary Hida realization of the weight-one theta form.
Shapiro therefore identifies the two base classes with the
ordinary Hida Beilinson--Flach specializations of KLZ
Definition~8.1.1, at the parameters in (CER.4).
This is equality of the defining universal class and its
quotient maps: the first map is the fixed moment and
ordinary-control map of Step~2, and the second is the actual
free rank-two CM specialization of Step~3.  In particular
we have not inferred equality of Galois cohomology classes
from equality on a dense set of weights, and we need no
comparison with a Coleman cohomology module at weight one.
The two choices of the CM type give the two ordinary theta
refinements.  This completes the construction and the base
identification.
\end{proof}

\begin{theorem}[AC-K29-lower: the remaining lower-boundary factor]
\label{AC-K29-lower}
Assume the hypotheses of AC-K29-ray, and suppose $f$ is
the non-CM ordinary type-II form of AC-K29-nonCM, with
the characteristic-zero image conclusion of AC-K29-split.
If
\[
                 L^{\rm imp}(f,\theta_\chi,e)\ne0,
                                                            \tag{CER.5}
\]
then
\[
             H^1_f(K,V\chi)=H^1_f(K,V\chi^2)=0.
                                                            \tag{CER.6}
\]
This is a coefficient-field vanishing statement.  It supplies
no uniform bound for the integral index of an Euler class.
\end{theorem}
\begin{proof}
Put $k=2e-1$, $k'=-1$, and $j=e-1$ in LZ (2016),
\S8.1, printed pp.~43--45.  These satisfy
$k'+1\le j\le k$.  The ordinary unit-root stabilization
of $f$ is noble.  The two roots for $g=\theta_\chi$ are
the distinct roots of unity $\omega,\omega^{-1}$; its
weight-one stabilizations are permitted by Theorem~4.7.2.
If $\alpha_f$ is the unit root and $\beta_f$ the other
root, their valuations are $0$ and $2e$.  Thus none of
$\alpha_f\alpha_g,\alpha_f\beta_g,
\beta_f\alpha_g,\beta_f\beta_g$ equals
$\ell^{e-1}$ or $\ell^e$: valuations exclude this when
$e>1$, and for $e=1$ the possible unit equality is excluded
by the complex absolute value $\ell^e$ of a root of $f$.
This verifies the source's exact no-local-zero condition.

We use the ordinary explicit reciprocity law directly,
namely KLZ Theorem~10.2.2, formula~(10.2.1), printed
pp.~52--53.  It is an identity for the universal ordinary
classes and their cuspidal quotient maps, before the
residual-image hypotheses imposed in the later Selmer
applications.  Its first-factor functional specializes at
our fixed positive-weight ordinary form by
Proposition~10.1.1(2)(b), and is a nonzero multiple of the
usual $f$-eigenform functional.  Its CM functional
$\omega_g$ is the canonical isomorphism of
Proposition~10.1.1(1); it is therefore nonzero on the actual
rank-one CM quotient at weight one.  The correction in
KLZ~\S13, printed p.~67, states that the corrected
congruence ideal localizes to the unit ideal at an arithmetic
point.  Thus no scalar vanishing at our fixed $f$ is
introduced by its reducible residual representation.

The required scalar interpolation is precisely KLZ
Theorem~7.7.2 and the sentence immediately following it,
printed p.~39: it extends to a classical weight-one
specialization of the second family.  Its formula is
Theorem~2.7.4; the same ordinary interpolation is given by
Loeffler, arXiv:1704.04049v2, Proposition~2.10 and
Theorem~6.3.  One may also identify the two scalar analytic
$p$-adic $L$-functions by their common interpolation on
classical critical triples in the full three-variable
families.  This compares analytic scalar functions, not
Galois cohomology classes.
At the full-weight point $(2e+1,1,e)$ the smoothing factor
in KLZ~(10.2.1) is $c^2-1$; choose a permitted $c$ with
$c^2\not\equiv1\pmod\ell$.  The interpolation Euler factors
are nonzero by the no-local-zero calculation above.
Thus (CER.5) makes the ordinary reciprocity scalar nonzero
for each of the two CM refinements.

We spell out why this scalar detects the actual singular
local class at weight one.  KLZ Theorem~8.2.3 constructs
the big logarithm functorially for every unramified
profinite module.  The projection of the universal class
to $F^{-\circ}$ lies in $F^{-+}$ by KLZ
Proposition~8.1.7, whose proof imposes no residual
irreducibility condition.  Apply the big logarithm to
$F^{-+}M(f\otimes g)^*(-1-k')$ and then specialize its
actual CM rank-one factor at $k'=-1$.
The module is still unramified; the twist $-1-k'$ is now
zero.  Its finite free specialization therefore gives the
usual Perrin--Riou map by the second assertion of that
theorem, with no algebraic moment of negative degree.
The proof of KLZ Theorem~8.2.8 is this same functorial
construction followed by a cyclotomic twist.  At
$j=e-1>k'=-1$ the resulting local map is
\[
 (e-1)!\,
 \frac{1-\ell^{e-1}/(\alpha_f\beta_g)}
      {1-\alpha_f\beta_g/\ell^e}\,\exp^*.
                                                            \tag{CER.5a}
\]
All scalar factors are nonzero.  The nonzero reciprocity
value consequently gives a nonzero singular image of each
base class.  Its $F^{--}$ projection is zero by
Proposition~8.1.7.  The CM representation splits locally
into its two unramified lines, and the ordinary filtration
of $f$ identifies the singular quotient with its
$F^-$ quotient.  Thus the two refinements give singular
images in the two distinct one-dimensional CM Frobenius
lines; together they span the two-dimensional singular
local space.  This is also the
local linear-algebra conclusion used in the proof of LZ
Theorem~8.1.4; the present argument obtains it on the
ordinary coefficient modules themselves.

Write $W=V\otimes\operatorname{Ind}_K^{\mathbb Q}\chi$.
The local part of the Poitou--Tate proof of LZ
Theorem~8.2.1, printed p.~46, says that the images of
\[
 H^1_f(\mathbb Q,W)\longrightarrow H^1_f(\mathbb Q_\ell,W)
 \quad\hbox{and}\quad
 H^1_{\mathrm{rel}(\ell)}(\mathbb Q,W^*(1))
       \longrightarrow H^1_s(\mathbb Q_\ell,W^*(1))
                                                            \tag{CER.7}
\]
are orthogonal complements.  The second image is the whole
space, by the two classes.  Consequently the first image is
zero.  This argument uses no rank-one image element for
$W$; indeed AC-K30 proves that such an element does not
exist.  Shapiro now gives
\[
 H^1_f(K,R_a)=H^1_{\mathrm{strict}\{\mathfrak p,
                    \overline{\mathfrak p}\}}(K,R_a).
                                                            \tag{CER.8}
\]

At least one base class of AC-K29-ray is nonzero, because
its singular localization is nonzero.  Apply the
characteristic-zero Euler-system bound in the proof of
LLZ Theorem~7.3.1, printed pp.~27--28, to $U_a$.
The lattice bound and away-$\ell$ local condition were
proved above.  Its image hypotheses are exactly the
irreducibility, rank-one quotient element and fixed-vector-free
element on the derived $G_K$ image furnished by
AC-K29-split, unchanged by the scalar twist in $U_a$.
Since all the auxiliary ray and cyclotomic extensions are
abelian, these derived-image elements act trivially on them.
The modified Rubin bound used there therefore shows that
the dual strict Selmer group over the coefficient field is
zero.  Together with (CER.8) this proves (CER.6).
No integral big-image theorem is used in this step.
\end{proof}

\begin{proposition}[AC-K29-bad: the exact shared-bad-factor gate]
\label{AC-K29-bad}
Retain AC-K29-lower's setting.  Let $M_g$ be the Artin
representation of $\theta_\chi$.  For a rational prime
$q\ne\ell$ put
\[
 C_q(X)=\det\left(1-X\operatorname{Fr}_q^{-1}\ \middle|\
 \frac{(M_f\otimes M_g)^{I_q}}
                  {M_f^{I_q}\otimes M_g^{I_q}}\right).
                                                            \tag{CER.9}
\]
Only primes shared by $N_f$ and $3\cdot269$ can contribute.
Then
\[
 L^{\rm imp}(f,g,e)
     =L^{\rm prim}(f,g,e)\prod_q C_q(q^{-e}).
                                                            \tag{CER.10}
\]
The primitive factor is nonzero.  Consequently the precise
remaining nonvanishing condition in (CER.5) is
$C_q(q^{-e})\ne0$ at every shared bad prime.

In the restricted subcase where $M_f$ is unramified at
$269$ and its conductor exponent at $3$ is one, this
condition is exactly
\[
                    a_3(f)\ne3^e.                         \tag{CER.11}
\]
Neither of these restricted local assumptions is imposed on
the full old family in AC-K29-point.
\end{proposition}
\begin{proof}
The denominator space in (CER.9) is an invariant subspace
of the numerator, so its determinant quotient is a
polynomial.  LZ (2016), the proof of Corollary~8.3.3, printed
p.~48, identifies the naive local Euler factor with
the determinant on $M_f^{I_q}\otimes M_g^{I_q}$.
Taking the determinant on the quotient gives (CER.10).
If either factor is unramified, the two invariant spaces
coincide, proving the asserted support of the product.

For the primitive value, the non-CM form has cuspidal
base change to $K$.  Normalize it to be unitary and twist
by $\chi$.  Jacquet--Shalika's nonvanishing theorem,
\emph{Inventiones}~38 (1976), Theorem~(1.3), printed p.~1,
gives a nonzero completed value at the line $\Re s=1$.
The functional equation gives a nonzero completed value
at $s=0$.  Here the complex archimedean factor is
$\Gamma_{\mathbb C}(s+e)^2$, which is finite and nonzero
at both $0$ and $1$ since $e\ge1$.  The finite primitive
value at zero is therefore nonzero.  Restoring the
classical weight shifts this point to $s=e$, as claimed.

For (CER.11), the Hensel calculation above shows
$M_g|_{G_{\mathbb Q_3}}=1\oplus\eta_K$.
Conductor exponent one and determinant $\eta_K$ imply
\[
       V|_{G_{\mathbb Q_3}}=A\oplus A^{-1}\eta_K
                                                            \tag{CER.12}
\]
with $A$ unramified.  Indeed the inertia-fixed line has
dimension one, Swan conductor zero, and the two distinct
tame inertia characters $1,\eta_K$ give invariant lines.
Their distinct finite inertia characters also force the
monodromy to vanish.  The determinant supplies the second
character in (CER.12).

Choose the Frobenius lift on which $\eta_K$ is one.
Such a lift exists because $3$ is the norm of $\sqrt{-3}$.
If $a=A(\operatorname{Fr}_3^{-1})$, then
$a=a_3(f)/3^e$.  The invariant line contributed by
$1\otimes1$ is the naive local factor; the additional
invariant line comes from $\eta_K\otimes\eta_K$ and
has normalized Frobenius eigenvalue $a^{-1}$.
Hence $C_3(3^{-e})=1-a^{-1}=1-3^e/a_3(f)$.
At $269$ the extra quotient is zero under the explicitly
assumed unramifiedness of $M_f$.  This proves (CER.11).
\end{proof}

\paragraph{What is still required for the whole source family.}
The old character $\psi$ may be ramified at $269$; the
old counting level contains $3$ and may contain $269$.
Thus (CER.11) is only the stated restricted calculation,
and a general member requires both relevant quotients
in (CER.9).  The coefficient prime also still needs the
ordinary good-level and cubic-distinction hypotheses of
AC-K29-ray, or a separately proved replacement.
If a factor in (CER.10) vanishes, dividing its numerical
value is not a construction of a cohomology class.
Primitive local test vectors would have to be realized by
compatible classes with the same norm and lattice properties.
Likewise, the right-boundary point $(2e+1,1,e+1)$ lies
on $k_1+k_2=2s$: the universal smoothing factor vanishes
there.  Dividing it in a meromorphic family would require
a proof of a regular cohomological specialization.
Neither repair is asserted here.

\paragraph{Realization and source scope.}
The construction takes place in the $\ell$-adic realizations
of modular-curve cohomology, using tensor products,
corestriction and the induced realization of the cubic
Artin character.  Its new assertion is the fixed-form
projection argument and the resulting scoped comparison;
the universal Rankin--Eisenstein classes, the CM integral
patching, the reciprocity law and Rubin's characteristic-zero
bound are cited inputs.  No new algebraic cycle over the
original prolate spaces, and no Grothendieck archival
identification, is claimed.
The relevant primary sources are
\url{https://arxiv.org/pdf/1503.02888v6},
\url{https://antoniolei.com/wp-content/uploads/2014/11/cmeuler.pdf},
\url{https://arxiv.org/pdf/1602.07508v3},
\url{https://link.springer.com/content/pdf/10.1186/s40687-016-0077-6.pdf},
\url{https://arxiv.org/pdf/1704.04049v2}, and
\url{https://doi.org/10.1007/BF01390166}.
